\origpage{46}{N. Robertson, P.D. Seymour / Journal of Combinatorial Theory, Series B 89 (2003) 43--76}
\textbf{1.3.} \textit{Let $L$ be a nonplanar graph, and let $\Sigma_1, \dots, \Sigma_s$ be all the connected surfaces (up to homeomorphism) in which $L$ cannot be drawn. Then there are numbers $r, d, w$ such that every graph with no $L$-minor may be constructed by clique-sums, starting from graphs $G'$ with the following property: $G'$ is a $(\le w)$-vertex extension of an outgrowth by $\le d$ $r$-rings of a graph that can be drawn in one of $\Sigma_1, \dots, \Sigma_s$.}

One can think of an $r$-ring as a graph with a given path-decomposition of bounded width, and in particular, its pieces come in a linear order. Since these pieces are being sewn onto a circular hole in a surface, it would perhaps be more natural if the pieces came in a circular order. One can indeed replace the third condition in the definition of an $r$-ring by the condition

\hspace*{12pt}$\bullet$ $X_i \cap X_k \subseteq X_h \cup X_j$ and $X_h \cap X_j \subseteq X_i \cup X_k$ for $1 \le h \le i \le j \le k \le n$

and 1.3 remains true (because the new statement is obviously weaker than the original); and in fact the new version is equivalent to the original (we leave the equivalence to the reader). This is why we call it a ``ring'' rather than some more linear name. It is a question of convenience which definition of ``$r$-ring'' is used.

While 1.3 has been one of the main goals of this series of papers, it turns out to have been a red herring. There is another result (theorem 3.1) which is proved in this paper, and from which 1.3 is then derived; and in all the future applications in this series of papers, it is not 1.3 but 3.1 that will be needed. Let us explain how theorem 3.1 is used to prove 1.3.

Evidently we would like to eliminate the ``tree-structure'' part of 1.3 and concentrate on the internal structure of one of the ``nodes'' of the tree. How can we do so? An inductive argument looks plausible at first sight; if there is no low order cutset of $G$ dividing it into two substantial pieces then $G$ itself must be almost a ``node'' if the theorem is to be true, while if there is such a cutset we may express $G$ as a clique-sum of two smaller graphs, and hope to apply our inductive hypothesis to these graphs. But there is a difficulty here; it is possible that these smaller graphs have an $L$-minor while $G$ does not. Fortunately there is a way to focus in on a ``node'' which does not involve any decomposing, as follows. We can assume that the tree is as refined as possible in the sense that no node can be split into two smaller nodes, and so for every low order cutset of $G$, most of any node will lie on one side or the other of the cutset (except for nodes of bounded cardinality, which we can ignore.) Therefore if we fix some node, every small cutset has a ``big'' side (containing most of the node) and a ``small'' side---and it turns out that no three small sides have union $G$. Thus a node defines a ``tangle'', which is such an assignment of big and small sides to the low order cutsets; and conversely, it can be shown that any tangle in $G$ of sufficiently high ``order'' will be associated with some node of the tree-structure. Hence a convenient way to analyze the internal structure of the nodes is to analyze the local structure of $G$ with respect to some high order tangle, and this is the content of theorem 3.1.
